\subsubsection{Cuantizador uniforme}
\label{sec:uniform}

El cuantizador uniforme consta de una serie de $n(b) = 2^b$ niveles equiespaciados dentro del rango $r(\alpha)$, donde $b$ es el número de bits empleados en la cuantización y $\alpha$ es un parámetro que define dicho rango.

El rango del cuantizador depende del tipo de cuantización:

\begin{itemize}
    \item En el caso signado, el rango es $r(\alpha) = 2\alpha$, abarcando el intervalo $[-\alpha, \alpha].$
    \item En el caso no signado, el rango es $r(\alpha) = \alpha$, abarcando el intervalo $[0, \alpha].$
\end{itemize}

De forma general, el intervalo del rango se define como \eqref{eq:range_interval}.

\begin{equation}
    \label{eq:range_interval}
    \mathbb{I} = [m, M]
\end{equation}


donde $m$ representa el valor mínimo y $M$ el valor máximo del intervalo del rango.

El paso de cuantización, denotado por $\Delta(\alpha, b)$, es la distancia entre niveles consecutivos y se define \eqref{eq:quantization_step}.

\begin{equation}
    \label{eq:quantization_step}
    \Delta(\alpha, b) = \frac{r(\alpha)}{n(b)}
\end{equation}

Los niveles de cuantización se expresan como \eqref{eq:uniform_quantizer_level_value}, donde $i$ representa el índice del nivel, que varía entre $0$ y $n(b)-1$.

\begin{equation}
    \label{eq:uniform_quantizer_level_value}
    l_i\left(\alpha; b\right) = m + i \cdot \Delta\left(\alpha, b\right) \text{, con } i = 0, 1, \ldots, n\left(b\right) - 1
\end{equation}

La función de cuantización uniforme se define \eqref{eq:uniform_quantizer} y su representación gráfica se muestra en la \reffig{uniform_signed}.

\begin{equation}
    \label{eq:uniform_quantizer}
    q_u(x; \alpha, b) =
    \begin{cases}
        l_{0} & \text{si $x \leq m$} \\
        \Delta \cdot \lfloor \frac{x}{\Delta} \rfloor & \text{si } m < x < M \\
        l_{n\left(b\right) - 1}\left(\alpha\right) & \text{si } x \geq M
    \end{cases}
\end{equation}

\defaultfigure{quantizers/uniform/q.png}{Función de cuantización uniforme signada.}{uniform_signed}

El cuantizador uniforme resulta sencillo de implementar tanto en hardware como en software debido a la distribución equiespaciada de sus niveles. Presenta eficiencia computacional, pues la cuantización se realiza mediante operaciones aritméticas simples, como divisiones y redondeos.

No obstante, existen escenarios en los que la cuantización uniforme no resulta adecuada.

Una aplicación donde los datos presentan una distribución no uniforme, esta puede no captar la variabilidad de los datos, lo que conduce a una pérdida significativa de información y a niveles en desuso. Si la aplicación exige alta precisión en determinadas regiones del rango, la cuantización uniforme puede no ofrecer la resolución necesaria en esas zonas críticas, este caso aparece con frecuencia en la distribución de los pesos de redes neuronales.

La cuantización uniforme resulta ineficiente al combinar rangos amplios con requisitos de alta precisión, puesto que la misma representación debe cubrir valores muy pequeños y muy grandes. Este enfoque puede ser viable con aritmética de punto flotante, no obstante, para aritmética entera o de punto fijo no es eficiente.

\subsubsection{Cuantizador no uniforme}
\label{sec:non_uniform}

El cuantizador no uniforme comprende una serie de $n$ niveles $\mathbf{l}$ cuyos límites están determinados por los umbrales $\mathbf{t}$. En este esquema no existen restricciones sobre la posición de umbrales y niveles, ni la noción de bits.

La función de cuantización no uniforme se define como \eqref{eq:non_uniform_quantizer} y su representación gráfica se muestra en la \reffig{non_uniform_signed_q}.

\begin{equation}
    \label{eq:non_uniform_quantizer}
    q_{nu}\left(x; \mathbf{t}, \mathbf{l}\right) =
    \begin{cases}
        l_{0} & \text{si $x \leq m$} \\
        l_{i} & \text{si } t_i < x < t_{i+1}, i = 0, 1, \ldots, n-2 \\
        l_{n-1} & \text{si } x \geq M
    \end{cases}
\end{equation}

donde $m$ y $M$ representan los límites del rango, es decir, $m = t_0$ y $M = t_{n-1}$.

\defaultfigure{quantizers/non_uniform/q.png}{Función de cuantización no uniforme signada.}{non_uniform_signed_q}

El cuantizador no uniforme es más versátil que el uniforme, pues permite adaptar los niveles a la distribución de los datos o a las necesidades del problema, lo que mejora la precisión y aumenta el nivel de compresión. No obstante, su implementación resulta más compleja, especialmente en hardware de punto fijo, debido a la incompatibilidad de algunos valores con la aritmética de punto fijo \cite{electronics13101923}.

% \TODO{Ver si hay que desarrollar xq es malo para punto fijo}

\subsubsection{Cuantizador flexible}
\label{sec:flex}

El cuantizador flexible pretende combinar la eficiencia de implementación del cuantizador uniforme con la capacidad del cuantizador no uniforme para adaptarse a distintas distribuciones de datos. Este método cuantiza uniformemente los niveles de un cuantizador no uniforme \cite{electronics13101923}.

El cuantizador flexible consta de un conjunto de $n$ niveles, con valores $\mathbf{l}$ cuyos límites están determinados por los umbrales $\mathbf{t}$. Como los casos anteriores, el rango se fija mediante el parámetro $\alpha$ y el rango del cuantizador depende del tipo de cuantización:

\begin{itemize}
    \item En el caso signado, el rango es $r(\alpha) = 2\alpha$, abarcando el intervalo $[-\alpha, \alpha].$
    \item En el caso no signado, el rango es $r(\alpha) = \alpha$, abarcando el intervalo $[0, \alpha].$
\end{itemize}

Para optimizar la implementación en la biblioteca, se definen $n$ niveles y $n+1$ umbrales. Esta elección puede resultar poco intuitiva, pero facilita la implementación posterior. En esta definición, los umbrales $t_0$ y $t_n$ son fijos y coinciden con los límites del rango: $t_0 = m$ y $t_n = M$.

La función de cuantización flexible se define \eqref{eq:flexible_quantizer}

\begin{equation}
    \label{eq:flexible_quantizer}
    q_{f}\left(x; \alpha, \mathbf{t}, \mathbf{l}\right) =
    \begin{cases}
        q_{u}\left(l_{0}; \alpha\right)& \text{si $x \leq m$} \\
        q_{u}\left(l_{i}; \alpha\right) & \text{si } t_i < x < t_{i+1}, i = 0, 1, \ldots, n-2 \\
        q_{u}\left(l_{n-1}; \alpha\right) & \text{si } x \geq M
    \end{cases}
\end{equation}

    donde $q_{u}$ es la función de cuantización uniforme definida \eqref{eq:uniform_quantizer}. El gráfico de la función de cuantización flexible se muestra en la \reffig{flexible_signed_q}.

\defaultfigure{quantizers/flex/q_q.png}{Función de cuantización flexible signada.}{flexible_signed_q}

La \reffig{flexible_signed_q} muestra, en amarillo, los niveles previos a la aplicación del cuantizador uniforme y, en azul, los niveles finales tras aplicar dicho cuantizador.

La principal diferencia respecto al uniforme es que los niveles $n$ no dependen de la cantidad de bits, tanto $\mathbf{l}$ como $\mathbf{t}$ son parámetros del cuantizador. En otras palabras, se desacoplan los niveles de cuantización de la cantidad de bits.

Este cuantizador aproxima la cuantización no uniforme mediante una función cuyos valores de salida son compatibles con la aritmética de punto fijo, por ello, permite una implementación eficiente en hardware en conjunto con la flexibilidad para adaptarse a distintas distribuciones de datos.
